% !TEX spellcheck = en_US



% ~~~~~~~~~~~~~~~~~~~~~~~~~~~~~~~~~~~~~~~ V
% (NC) 2026 Haluk Bingol
% github.com/halukbingol/LaTeX-Templates
%
% Licensees may copy, distribute, display, and perform the work and make derivative works 
% and remixes based on it only for non-commercial purposes. 
% ~~~~~~~~~~~~~~~~~~~~~~~~~~~~~~~~~~~~~~~ A




% =====================================================================
%  Maven for Java -- Freshman Level: Install, First Project, Build and Test
%  ---------------------------------------------------------------------
%  Built on hblib/Template/hblibTemplate.tex with the libraries in ../../hblib/.
%  Programs and their real outputs are in codeoutput/:
%     codeoutput/<Name>.java  a program (classes without main() are only compiled)
%     codeoutput/<Name>.sh    a terminal session, run in a copy of codeoutput/<Name>/
%     codeoutput/<Name>.txt   what it printed, made by:  bash run_examples.sh
% =====================================================================

\documentclass[aspectratio=169,10pt,compress]{beamer}

% ~~~~~~~~~~~~~~~~~~~~~~~~~~~~~~~~~~~~~~~~~~~~~~~~~~~~~~~~~~~~~~~~~ V hblib
	\usepackage{../../hblib/hblibPresentation} % lib presentation: theme, i/N, \hSection, algorithm2e, ...
	\usepackage{../../hblib/hblibCodeoutput} % lib code + output: \lstcode, \lstterminal, ...
	\usepackage{../../hblib/hblibDatastructures} % lib data structures: \hString, \hStack, \hQueue
% ~~~~~~~~~~~~~~~~~~~~~~~~~~~~~~~~~~~~~~~~~~~~~~~~~~~~~~~~~~~~~~~~~ A hblib


% theme, i/N, fonts (T1, lmodern), booktabs, xcolor, array (L{width} columns), listings, tikz,
% algorithm2e (\small captions, \Fn) and the color gitred come from the hblib libraries




% +++++++++++++++++++++++++++++++++++++++ V specific
\DontPrintSemicolon % in algorithms
\newcommand{\win}{\textbf{Windows}}
\newcommand{\mac}{\textbf{macOS}}
% +++++++++++++++++++++++++++++++++++++++ A specific

% ---------------------------------------------------------------------
\title{Maven for Java}
\subtitle{Freshman Level: Install, First Project, Build and Test}
\author[Bingol]{Haluk O. Bingol}

\institute[FCIS]{
    Faculty of Computer and Information Sciences\\
    Yeditepe University
}
\date{
	{\tiny \hbTimeStamp}
}

\begin{document}




% ~~~~~~~~~~~~~~~~~~~~~~~~~~~~~~~~~~~~~~ V frm
\begin{frame}[t,fragile]
  \titlepage
\end{frame}
% ~~~~~~~~~~~~~~~~~~~~~~~~~~~~~~~~~~~~~~ A frm




% ~~~~~~~~~~~~~~~~~~~~~~~~~~~~~~~~~~~~~~ V frm
\begin{frame}[t,fragile] \frametitle{Outline}
  \begin{columns}[T]
    \begin{column}{0.48\textwidth}
      \tableofcontents[sections={1-4}, hideallsubsections]
    \end{column}
    \begin{column}{0.48\textwidth}
      \tableofcontents[sections={5-7}, hideallsubsections]
    \end{column}
  \end{columns}
\end{frame}
% ~~~~~~~~~~~~~~~~~~~~~~~~~~~~~~~~~~~~~~ A frm

% ~~~~~~~~~~~~~~~~~~~~~~~~~~~~~~~~~~~~~~ V
\hSection[Why]{Why a Build Tool?}
% ~~~~~~~~~~~~~~~~~~~~~~~~~~~~~~~~~~~~~~ A




% ~~~~~~~~~~~~~~~~~~~~~~~~~~~~~~~~~~~~~~ V frm
\begin{frame}[t,fragile] \frametitle{Building by hand does not scale}
  \begin{columns}[T]
    \begin{column}{0.48\textwidth}
      \textbf{By hand}
      \begin{itemize}\small
        \item \texttt{javac -d out -cp lib/a.jar:lib/b.jar \dots} for every file.
        \item Download each library (\texttt{.jar}) and the libraries \emph{it} needs.
        \item Run the tests yourself, put the classes in a \texttt{.jar} yourself.
        \item Every student and every IDE does it a little differently.
      \end{itemize}
    \end{column}
    \begin{column}{0.48\textwidth}
      \textbf{With Maven}
      \begin{itemize}\small
        \item One file, \texttt{pom.xml}, describes the project.
        \item \texttt{mvn package}: compile, test, make the \texttt{.jar}.
        \item Libraries are named, Maven downloads them (and theirs).
        \item The same build on Windows, macOS, Linux, in the IDE and on a build server.
      \end{itemize}
    \end{column}
  \end{columns}
  \vspace{0.5em}
  \hRemark{Maven: \emph{convention over configuration}: follow the standard layout and you write almost nothing.}
\end{frame}
% ~~~~~~~~~~~~~~~~~~~~~~~~~~~~~~~~~~~~~~ A frm

% ~~~~~~~~~~~~~~~~~~~~~~~~~~~~~~~~~~~~~~ V
\hSection[Install]{Installing Maven}
% ~~~~~~~~~~~~~~~~~~~~~~~~~~~~~~~~~~~~~~ A




% ~~~~~~~~~~~~~~~~~~~~~~~~~~~~~~~~~~~~~~ V frm
\begin{frame}[t,fragile] \frametitle{Getting Maven}
  \begin{itemize}\small
    \item Needs a JDK first (\texttt{java -version}, \texttt{javac -version}); here JDK 21.
    \item Current release: Maven 3.10.0 (September 2026); Maven 4 is still a release candidate. The sessions in this deck ran Maven 3.9.11: the commands are the same.
  \end{itemize}
  \begin{columns}[T]
    \begin{column}{0.48\textwidth}
      \begin{block}{Windows (not run here)}


% +++++++++++++++++++++++++++++++++++++++ V lst
%: -Lst.
\begin{lstlisting}[style=codesnippet, language={}, basicstyle=\ttfamily\tiny]
:: 1. maven.apache.org: download
::    apache-maven-3.10.0-bin.zip
:: 2. unzip to C:\tools\apache-maven-3.10.0
:: 3. Settings > System > About >
::    Advanced system settings >
::    Environment Variables: add to Path
::    C:\tools\apache-maven-3.10.0\bin
:: 4. a NEW Command Prompt:
mvn -v
\end{lstlisting}
% +++++++++++++++++++++++++++++++++++++++ A lst


      \end{block}
    \end{column}
    \begin{column}{0.48\textwidth}
      \begin{block}{macOS (not run here)}


% +++++++++++++++++++++++++++++++++++++++ V lst
%: -Lst.
\begin{lstlisting}[style=codesnippet, language=bash, basicstyle=\ttfamily\tiny]
# with Homebrew (brew.sh):
brew install maven
mvn -v

# or by hand: download the .tar.gz,
# unpack it, then in ~/.zshrc:
export PATH="$HOME/tools/apache-maven-3.10.0/bin:$PATH"
\end{lstlisting}
% +++++++++++++++++++++++++++++++++++++++ A lst


      \end{block}
    \end{column}
  \end{columns}

  \small If \texttt{mvn} says \texttt{JAVA\_HOME} is not defined: set \texttt{JAVA\_HOME} to the JDK folder the same way.
\end{frame}
% ~~~~~~~~~~~~~~~~~~~~~~~~~~~~~~~~~~~~~~ A frm




% ~~~~~~~~~~~~~~~~~~~~~~~~~~~~~~~~~~~~~~ V frm
\begin{frame}[t,fragile] \frametitle{Is it working?}
  \lstterminal[basicstyle=\ttfamily\scriptsize]{codeoutput/Version.txt}

  \vspace{0.1em}
  \small
  \begin{itemize}
    \item The Maven version, where it is installed, and the Java it uses: check that this is the JDK you expect.
    \item Run here on Linux; on Windows and macOS the lines look the same, with other paths.
  \end{itemize}
\end{frame}
% ~~~~~~~~~~~~~~~~~~~~~~~~~~~~~~~~~~~~~~ A frm

% ~~~~~~~~~~~~~~~~~~~~~~~~~~~~~~~~~~~~~~ V
\hSection[Project]{The First Project}
% ~~~~~~~~~~~~~~~~~~~~~~~~~~~~~~~~~~~~~~ A




% ~~~~~~~~~~~~~~~~~~~~~~~~~~~~~~~~~~~~~~ V frm
\begin{frame}[t,fragile] \frametitle{The standard layout}
  \begin{columns}[T]
    \begin{column}{0.56\textwidth}
      \lstcode[basicstyle=\ttfamily\tiny]{codeoutput/Layout/src/main/java/hello/Greeter.java}{Java}
    \end{column}
    \begin{column}{0.42\textwidth}
      \uncover<2->{\lstterminal[basicstyle=\ttfamily\scriptsize]{codeoutput/Layout.txt}}
    \end{column}
  \end{columns}

  \vspace{0.1em}
  \small
  Sources in \texttt{src/main/java}, tests in \texttt{src/test/java}, folders = packages. (\texttt{tree}: a small shell function here; Windows: \texttt{tree /F}.)
\end{frame}
% ~~~~~~~~~~~~~~~~~~~~~~~~~~~~~~~~~~~~~~ A frm




% ~~~~~~~~~~~~~~~~~~~~~~~~~~~~~~~~~~~~~~ V frm
\begin{frame}[t,fragile] \frametitle{\texttt{pom.xml}: the Project Object Model}
  \begin{columns}[T]
    \begin{column}{0.5\textwidth}
      \lstcode[basicstyle=\ttfamily\tiny, lastline=23]{codeoutput/Layout/pom.xml}{XML}
    \end{column}
    \begin{column}{0.48\textwidth}
      \begin{itemize}\small
        \item \textbf{Coordinates}: \texttt{groupId} (who, reversed domain), \texttt{artifactId} (what), \texttt{version}.
        \item \texttt{SNAPSHOT}: still in development.
        \item \texttt{maven.compiler.release}: compile for Java 21.
        \item \texttt{UTF-8}: Turkish letters in the sources.
        \item One dependency: JUnit, only for the tests (\texttt{scope test}).
      \end{itemize}
    \end{column}
  \end{columns}
\end{frame}
% ~~~~~~~~~~~~~~~~~~~~~~~~~~~~~~~~~~~~~~ A frm




% ~~~~~~~~~~~~~~~~~~~~~~~~~~~~~~~~~~~~~~ V frm
\begin{frame}[t,fragile] \frametitle{\texttt{pom.xml}: pinning the plugin versions}
  \begin{columns}[T]
    \begin{column}{0.62\textwidth}
      \lstcode[basicstyle=\ttfamily\tiny, firstline=24, firstnumber=24]{codeoutput/Layout/pom.xml}{XML}
    \end{column}
    \begin{column}{0.36\textwidth}
      \begin{itemize}\small
        \item Maven itself does little: \textbf{plugins} compile, test, pack.
        \item Without versions here, Maven picks its defaults: the build may change when Maven changes.
        \item The versions here are those available offline for this deck; newer ones exist.
      \end{itemize}
    \end{column}
  \end{columns}
\end{frame}
% ~~~~~~~~~~~~~~~~~~~~~~~~~~~~~~~~~~~~~~ A frm

% ~~~~~~~~~~~~~~~~~~~~~~~~~~~~~~~~~~~~~~ V
\hSection[Build]{Compile, Test, Package}
% ~~~~~~~~~~~~~~~~~~~~~~~~~~~~~~~~~~~~~~ A




% ~~~~~~~~~~~~~~~~~~~~~~~~~~~~~~~~~~~~~~ V frm
\begin{frame}[t,fragile] \frametitle{\texttt{mvn compile}}
  \lstterminal[basicstyle=\ttfamily\tiny]{codeoutput/Compile.txt}

  \vspace{0.1em}
  \small
  The classes go to \texttt{target/classes}; run them with \texttt{java -cp} (the same on Windows and macOS).
\end{frame}
% ~~~~~~~~~~~~~~~~~~~~~~~~~~~~~~~~~~~~~~ A frm




% ~~~~~~~~~~~~~~~~~~~~~~~~~~~~~~~~~~~~~~ V frm
\begin{frame}[t,fragile] \frametitle{\texttt{mvn test} (1/2)}
  \lstterminal[basicstyle=\ttfamily\tiny, lastline=21]{codeoutput/Test.txt}

  \vspace{0.1em}
  \small Each \texttt{---} line is one plugin goal at work.
\end{frame}
% ~~~~~~~~~~~~~~~~~~~~~~~~~~~~~~~~~~~~~~ A frm




% ~~~~~~~~~~~~~~~~~~~~~~~~~~~~~~~~~~~~~~ V frm
\begin{frame}[t,fragile] \frametitle{\texttt{mvn test} (2/2)}
  \lstterminal[basicstyle=\ttfamily\tiny, firstline=22]{codeoutput/Test.txt}

  \vspace{0.1em}
  \small Surefire runs the JUnit tests: 1 run, 0 failures. The time and date change with every run.
\end{frame}
% ~~~~~~~~~~~~~~~~~~~~~~~~~~~~~~~~~~~~~~ A frm




% ~~~~~~~~~~~~~~~~~~~~~~~~~~~~~~~~~~~~~~ V frm
\begin{frame}[t,fragile] \frametitle{Reading the log}
  \begin{itemize}\small
    \item \texttt{--- compiler:3.10.1:compile (default-compile) @ hello ---}: plugin \texttt{compiler}, version, \textbf{goal} \texttt{compile}, project \texttt{hello}.
    \item \texttt{mvn test} ran everything before it: resources, compile, test resources, test compile, then the tests.
    \item These steps are \textbf{phases} of the \textbf{lifecycle}:
  \end{itemize}
  \[
  	\mathit{validate} \to \mathit{compile} \to \mathit{test} \to \mathit{package} \to \mathit{verify} \to \mathit{install} \to \mathit{deploy}
  \]
  \begin{itemize}\small
    \item Naming a phase runs all phases up to it. \texttt{BUILD SUCCESS} or \texttt{BUILD FAILURE} is the verdict.
  \end{itemize}
\end{frame}
% ~~~~~~~~~~~~~~~~~~~~~~~~~~~~~~~~~~~~~~ A frm




% ~~~~~~~~~~~~~~~~~~~~~~~~~~~~~~~~~~~~~~ V frm
\begin{frame}[t,fragile] \frametitle{\texttt{mvn package} and \texttt{mvn clean}}
  \lstterminal[basicstyle=\ttfamily\scriptsize]{codeoutput/Package.txt}

  \vspace{0.1em}
  \small
  \texttt{-q} (quiet): only errors. The \texttt{.jar} name is \texttt{artifactId-version.jar}. \texttt{clean} deletes \texttt{target}: everything there can be rebuilt.
\end{frame}
% ~~~~~~~~~~~~~~~~~~~~~~~~~~~~~~~~~~~~~~ A frm




% ~~~~~~~~~~~~~~~~~~~~~~~~~~~~~~~~~~~~~~ V frm
\begin{frame}[t,fragile] \frametitle{\texttt{java -jar} does not work (yet)}
  \lstterminal[basicstyle=\ttfamily\scriptsize]{codeoutput/JarRun.txt}

  \vspace{0.1em}
  \small
  The jar does not say which class has \texttt{main}. Either name the class (\texttt{java -cp \dots\ Greeter}, previous slide) or add a \emph{manifest} to the jar plugin (sophomore level).
\end{frame}
% ~~~~~~~~~~~~~~~~~~~~~~~~~~~~~~~~~~~~~~ A frm

% ~~~~~~~~~~~~~~~~~~~~~~~~~~~~~~~~~~~~~~ V
\hSection[Errors]{When the Build Fails}
% ~~~~~~~~~~~~~~~~~~~~~~~~~~~~~~~~~~~~~~ A




% ~~~~~~~~~~~~~~~~~~~~~~~~~~~~~~~~~~~~~~ V frm
\begin{frame}[t,fragile] \frametitle{A compile error}
  \lstterminal[basicstyle=\ttfamily\tiny]{codeoutput/CompileError.txt}

  \vspace{0.1em}
  \small
  \texttt{[5,40]}: line 5, column 40 of \texttt{Greeter.java}. Read the \emph{first} error; the \texttt{[Help 1]} lines are the same for every failure.
\end{frame}
% ~~~~~~~~~~~~~~~~~~~~~~~~~~~~~~~~~~~~~~ A frm




% ~~~~~~~~~~~~~~~~~~~~~~~~~~~~~~~~~~~~~~ V frm
\begin{frame}[t,fragile] \frametitle{A failing test}
  \lstterminal[basicstyle=\ttfamily\tiny]{codeoutput/FailTest.txt}

  \vspace{0.1em}
  \small
  \texttt{expected} vs.\ \texttt{but was}: the comma is missing. No \texttt{.jar} is made while a test fails.
\end{frame}
% ~~~~~~~~~~~~~~~~~~~~~~~~~~~~~~~~~~~~~~ A frm




% ~~~~~~~~~~~~~~~~~~~~~~~~~~~~~~~~~~~~~~ V frm
\begin{frame}[t,fragile] \frametitle{No \texttt{pom.xml} here}
  \lstterminal[basicstyle=\ttfamily\tiny]{codeoutput/NoPom.txt}

  \vspace{0.1em}
  \small
  Maven looks for \texttt{pom.xml} in the current folder: \texttt{cd} into the project first (Windows: \texttt{cd hello}, \texttt{dir}).
\end{frame}
% ~~~~~~~~~~~~~~~~~~~~~~~~~~~~~~~~~~~~~~ A frm




% ~~~~~~~~~~~~~~~~~~~~~~~~~~~~~~~~~~~~~~ V frm
\begin{frame}[t,fragile] \frametitle{A typo in the phase}
  \lstterminal[basicstyle=\ttfamily\tiny]{codeoutput/BadPhase.txt}

  \vspace{0.1em}
  \small
  \texttt{complie} is not a phase; the message lists all valid ones.
\end{frame}
% ~~~~~~~~~~~~~~~~~~~~~~~~~~~~~~~~~~~~~~ A frm




% ~~~~~~~~~~~~~~~~~~~~~~~~~~~~~~~~~~~~~~ V frm
\begin{frame}[t,fragile] \frametitle{A typo in the XML}
  \lstterminal[basicstyle=\ttfamily\tiny]{codeoutput/BadXml.txt}

  \vspace{0.1em}
  \small
  \texttt{</versoin>}: the POM is not valid XML; line 19, column 32 is given. An IDE marks such errors while you type.
\end{frame}
% ~~~~~~~~~~~~~~~~~~~~~~~~~~~~~~~~~~~~~~ A frm




% ~~~~~~~~~~~~~~~~~~~~~~~~~~~~~~~~~~~~~~ V frm
\begin{frame}[t,fragile] \frametitle{The test library is missing}
  \lstterminal[basicstyle=\ttfamily\tiny]{codeoutput/NoJunit.txt}

  \vspace{0.1em}
  \small
  Without the JUnit dependency the test does not compile: \texttt{package \dots\ does not exist}. (\texttt{| head -12}: the first 12 lines; the rest repeats them.)
\end{frame}
% ~~~~~~~~~~~~~~~~~~~~~~~~~~~~~~~~~~~~~~ A frm

% ~~~~~~~~~~~~~~~~~~~~~~~~~~~~~~~~~~~~~~ V
\hSection[IDE]{Maven in the IDE}
% ~~~~~~~~~~~~~~~~~~~~~~~~~~~~~~~~~~~~~~ A




% ~~~~~~~~~~~~~~~~~~~~~~~~~~~~~~~~~~~~~~ V frm
\begin{frame}[t,fragile] \frametitle{Opening a Maven project in an IDE}
  \begin{itemize}\small
    \item \textbf{Eclipse}: \menu{File > Import > Maven > Existing Maven Projects}, choose the folder with \texttt{pom.xml}. Run: right click \menu{Run As > Maven build\dots}, goal \texttt{package}.
    \item \textbf{IntelliJ IDEA}: \menu{File > Open}, choose \texttt{pom.xml}, \emph{Open as Project}; the \emph{Maven} tool window lists the phases.
    \item \textbf{VS Code}: the \emph{Extension Pack for Java} reads \texttt{pom.xml} on opening the folder.
    \item The IDE uses the POM: changing a dependency in \texttt{pom.xml} changes the IDE's class path.
    \item Same on Windows and macOS; \texttt{mvn} in a terminal stays the reference: if it builds there, it builds everywhere.
  \end{itemize}
\end{frame}
% ~~~~~~~~~~~~~~~~~~~~~~~~~~~~~~~~~~~~~~ A frm

% ~~~~~~~~~~~~~~~~~~~~~~~~~~~~~~~~~~~~~~ V
\hSection[Summary]{Summary}
% ~~~~~~~~~~~~~~~~~~~~~~~~~~~~~~~~~~~~~~ A




% ~~~~~~~~~~~~~~~~~~~~~~~~~~~~~~~~~~~~~~ V frm
\begin{frame}[t,fragile] \frametitle{Summary}
  \begin{center}\small
  \begin{tabular}{@{}L{0.32\textwidth}L{0.6\textwidth}@{}}
    \toprule
    \textbf{command}            & \textbf{what it does} \\
    \midrule
    \texttt{mvn -v}             & Maven and Java versions \\
    \texttt{mvn compile}        & \texttt{src/main/java} $\to$ \texttt{target/classes} \\
    \texttt{mvn test}           & compile, then run the tests in \texttt{src/test/java} \\
    \texttt{mvn package}        & \dots then \texttt{target/artifactId-version.jar} \\
    \texttt{mvn clean}          & delete \texttt{target} \\
    \texttt{mvn clean package}  & both, in this order \\
    \texttt{-q}                 & quiet: only errors \\
    \bottomrule
  \end{tabular}
  \end{center}
  \begin{itemize}\small
    \item The project is described by its coordinates and its \texttt{pom.xml}; layout by convention.
    \item On failure: read the first \texttt{[ERROR]}, it names the file, line and column.
  \end{itemize}
\end{frame}
% ~~~~~~~~~~~~~~~~~~~~~~~~~~~~~~~~~~~~~~ A frm




% ~~~~~~~~~~~~~~~~~~~~~~~~~~~~~~~~~~~~~~ V frm
\begin{frame}[t,fragile] \frametitle{Exercises}
  \begin{enumerate}\small
    \item Install Maven on your computer; compare your \texttt{mvn -v} with the one in the deck.
    \item Rebuild the \texttt{hello} project by hand (folders, \texttt{pom.xml}, two classes); run \texttt{mvn package}.
    \item Add a method \texttt{farewell(name)} returning \texttt{"Hoşça kal, name!"} and a test for it.
    \item Break the build in three ways (syntax, test, POM) and explain each first \texttt{[ERROR]} line.
    \item Change \texttt{version} to \texttt{1.0}; what is the name of the jar now?
  \end{enumerate}
\end{frame}
% ~~~~~~~~~~~~~~~~~~~~~~~~~~~~~~~~~~~~~~ A frm




% ~~~~~~~~~~~~~~~~~~~~~~~~~~~~~~~~~~~~~~ V frm
\begin{frame}[t,fragile]
	\begin{center}
		\vspace{4cm}
		\hRemark{\Huge Questions?}
	\end{center}
\end{frame}
% ~~~~~~~~~~~~~~~~~~~~~~~~~~~~~~~~~~~~~~ A frm



\end{document}
